%% toc: P15 173 173.5,173.6,173.7,173.8
\chapter{A Range Estimate and Its Omitted Variable}
\label{sec:p15-w02b-p15-073}

An order-of-magnitude estimate is a calculation inside a model. Giving ranges
for its inputs makes the calculation honest about the inputs it lists; it
says nothing about the quantities the model leaves out. This chapter works
one estimate from start to finish: the number of people a single highway
lane can carry per hour. It gives the estimate with explicit input ranges,
reports only the digits those ranges support, and then exhibits a variable
the estimate omits, vehicle length divided by speed, together with the
regime in which that omission makes the stated range wrong by more than an
order of magnitude.

The inputs are a model function, a stipulated range for each listed input,
and a conclusion to be tested, here a threshold of a thousand people per hour.
The outputs are the exact image of the input ranges, a rounded report of it,
a verdict on the threshold, and a description of where the estimate stops
applying. The dimension-matrix calculation of
\cref{thm:p15-dimensional-rank-nullity} identifies the dimensionless group
that controls the omission. The product-range bounds of
\cref{thm:p15-dimensional-product-bounds} and the decision certificates of
\cref{thm:p15-approximation-certificates} are used as stated there.

\section{The estimate and its source}
\label{sec:p15-w02b-p15-073-source}

Mahajan compares the rush-hour capacity of rail and road, defining capacity
as passengers transported per unit time. For one highway lane he uses the
two-second following rule: each driver leaves two seconds of travel time
behind the car in front, so that the lane carries one car every two seconds.
With roughly one person per car, the capacity is one person every
\(2\,\mathrm s\), that is \(1800\) people per hour
\cite[p.~6, eq.~(1.3)]{Mahajan2014Insight}.
He immediately qualifies the value: the two-second rule is not a law of
nature, and the average gap might be about \(4\,\mathrm s\) late at night and
\(1\,\mathrm s\) during the day; the single value is a compromise
\cite[p.~7]{Mahajan2014Insight}. He sets this lane capacity beside a
train-line estimate of \(18\,000\) people per hour obtained from
\(150\) people per car, \(20\) cars per train and \(6\) trains per hour
\cite[p.~7, eq.~(1.4)]{Mahajan2014Insight}.

The source's lane calculation reads ``one car every two seconds'' directly
from the following rule. The rule fixes the travel time between the rear of
one car and the front of the next. The time between successive cars passing
a fixed point also includes the time a car takes to pass its own length.
That term is the omitted variable studied below. It is not an error in the
source's purpose, which was a rough comparison at highway speed; it is the
limit of the estimate when it is carried to other speeds.

\section{Definitions and assumptions}
\label{sec:p15-w02b-p15-073-definitions}

\begin{hbdefinition}{def:p15-w02b-p15-073-range-estimate}{Range estimate}
A \emph{range estimate} consists of a model function
\(f:B\to(0,\infty)\) and an \emph{input box}
\(B=\prod_{i=1}^k[l_i,u_i]\) with \(0<l_i\leq u_i\).
Each interval \([l_i,u_i]\) is a stipulated range for the \(i\)th listed
input. The \emph{range} of the estimate is the image \(f(B)\).
The estimate \emph{supports} a statement \(S(y)\) about the target \(y\)
when \(S(f(x))\) holds for every \(x\in B\).
A range estimate is conditional on two premises: that the target equals
\(f\) at the actual inputs, and that the actual inputs lie in \(B\).
\end{hbdefinition}

\begin{hbdefinition}{def:p15-w02b-p15-073-omitted-variable}{Omitted variable of an estimate}
Let \(f\) be the model of a range estimate. A quantity \(w\) not among the
inputs of \(f\) is an \emph{omitted variable} relative to a fuller model
\(F(x,w)\) when the fuller model reduces to \(f\) in a limit or at a reference
value \(w_0\) of \(w\), that is \(F(x,w)\to f(x)\) as \(w\to w_0\).
The omission \emph{fails} a supported statement when some admissible
\(w\) gives \(F(x,w)\) values for which the statement is false.
The omission \emph{falsifies the range} when some admissible \((x,w)\)
gives \(F(x,w)\notin f(B)\).
\end{hbdefinition}

The fuller model is itself a premise. Calling \(w\) omitted asserts only that
\(F\) contains it and \(f\) does not; whether \(F\) describes the world is a
separate empirical question.

\begin{hbconvention}{def:p15-w02b-p15-073-reporting}{Reporting a range without unearned digits}
State the input box and the exact image \(f(B)\) so that the arithmetic can
be checked. For the conclusion, round the lower endpoint \emph{down} and the
upper endpoint \emph{up}, each to one significant figure; the rounded
interval still contains \(f(B)\). Give an order of magnitude \(10^j\) only
when the rounded interval lies within \([10^j,10^{j+1}]\); otherwise report
the rounded interval alone. Do not report a central value, such as an
arithmetic or geometric midpoint, with more digits than the endpoints
support: the input box assigns no preferred point inside the range.
\end{hbconvention}

\begin{hbdefinition}{def:p15-w02b-p15-073-lane-model}{Steady single-lane traffic}
In \emph{steady single-lane traffic} every car has length \(L>0\), moves at
the same constant speed \(v>0\), and carries \(n>0\) people. The front of
each car is a fixed distance \(s\geq L\), the \emph{spacing}, behind the front
of the car ahead. The \emph{car flow} is the number of car fronts crossing a
fixed point per unit time, and the \emph{lane flow} \(q\) is \(n\) times the
car flow, in people per unit time.

In the \emph{time-gap model} each driver keeps a travel time \(\tau>0\)
between the rear of the car ahead and the driver's own front, so that
\(s=L+v\tau\). The \emph{following-rule estimate} is
\[
 q_0=\frac n\tau .
\]
Here \(n\) is a count ratio and has no physical dimension; \(\tau\) is a
time, \(L\) a length and \(v\) a speed.
\end{hbdefinition}

The requirement \(s\geq L\) states only that cars do not overlap. The
time-gap model adds a behavioural premise, and the following-rule estimate
adds a further one, that the length term is negligible. All three are
model premises, not reported observations.

\section{Dimensional form of the omission}
\label{sec:p15-w02b-p15-073-dimensions}

Take the quantities \(q,\tau,L,v\) with \(n\) treated as a dimensionless
count. In the order (length, time) the dimension matrix, with columns
\(q,\tau,L,v\), is
\[
 D=\begin{pmatrix}0&0&1&1\\-1&1&0&-1\end{pmatrix}.
\]
It has rank two, so by \cref{thm:p15-dimensional-rank-nullity} there are
\(4-2=2\) independent dimensionless products. The kernel vectors
\((1,1,0,0)\) and \((0,-1,1,-1)\) give
\[
 \pi_1=q\tau,\qquad z=\frac{L}{v\tau}.
\]
A single-valued relation among these quantities, with \(n\) as a parameter,
therefore takes the form \(q\tau/n=\Phi(z)\) for some function \(\Phi\).
Dimensions alone leave \(\Phi\) undetermined. The following-rule estimate is
the assumption \(\Phi\equiv1\): it keeps the first product and discards the
second. The omitted variable enters only through \(z\), which compares the
time a car takes to pass its own length with the following gap.
\Cref{thm:p15-w02b-p15-073-lane-flow} shows that the time-gap model supplies
\(\Phi(z)=1/(1+z)\).

\section{Detecting a dimensionally invalid candidate}
\label{sec:p15-w02b-p15-073-invalid}

Before a candidate formula for \(q\) is given input ranges it should pass a
dimensional check. The check compares dimension vectors in the coordinates
of \cref{def:p15-dimensional-coordinate}, here in the order (length, time):
\[
 [n]=(0,0),\quad [\tau]=(0,1),\quad [L]=(1,0),\quad
 [v]=(1,-1),\quad [q]=(0,-1).
\]
Write a unit change as \(x_i\mapsto x_i e^{d_i^{\mathsf T}\sigma}\) with
\(\sigma\in\mathbb R^k\); \(\sigma\) plays the role of the vector \(s\) of
\cref{def:p15-dimensional-coordinate}, renamed because \(s\) is the spacing
in this chapter. A candidate is \emph{dimensionally valid} for a target of
dimension vector \(d_y\) when it transforms by the factor
\(e^{d_y^{\mathsf T}\sigma}\) under every unit change, as required in
\cref{thm:p15-dimensional-covariant-output}.

\begin{hblemma}{thm:p15-w02b-p15-073-invalid-candidate}{proved=annex:p15-w02b-p15-073-invalid-candidate}{Two tests that expose an invalid candidate}
Let positive quantities \(x_1,\ldots,x_m\) have dimension vectors
\(d_1,\ldots,d_m\in\mathbb R^k\), forming the matrix \(D\), and let the
target \(y>0\) have dimension vector \(d_y\).
\begin{enumerate}
\item \emph{Monomial test.} Let \(g(x)=Cx^a\) with a constant \(C>0\) that
does not depend on the units. If \(Da\neq d_y\) and \(y=g(x)\) in one system
of units, then after the unit change \(\sigma=t(d_y-Da)\), \(t\neq0\), the
transformed values satisfy
\[
 \frac{y'}{g(x')}=e^{t\lVert d_y-Da\rVert^2}\neq1 .
\]
\item \emph{Sum test.} If \(d_i\neq d_j\), then for every positive
\(x_i,x_j\) and every \(c>0\) some unit change gives
\(x_i'/x_j'=c\). Which term of \(x_i+x_j\) dominates is then fixed by the
choice of units, not by the quantities.
\item \emph{The candidate \(\tilde g=n/(\tau+L)\).} With the vectors
displayed above and \(\sigma=(\sigma_L,\sigma_T)\),
\(\tilde g(x')=e^{d_y^{\mathsf T}\sigma}\tilde g(x)\) for \(d_y=[q]\) holds
exactly when \(\sigma_L=\sigma_T\). With the time unit held fixed
\((\sigma_T=0)\), as the length unit varies over all positive multiples of
the original one, \(\tilde g\) takes every value in \((0,n/\tau)\), while
the target \(q\) does not change.
\end{enumerate}
\end{hblemma}

\begin{hbexample}{def:p15-w02b-p15-073-invalid-check}{Three candidate lane formulas checked}
Fix one physical state: \(n=1\), \(\tau=2\,\mathrm s\), \(L=5\,\mathrm m\),
\(v=30\,\mathrm{m/s}\). Each candidate below claims to give \(q\).
\begin{enumerate}
\item \(q\overset{?}{=}n/(\tau+L)\). The sum adds \([\tau]=(0,1)\) to
\([L]=(1,0)\); the vectors differ, so the candidate is invalid by
\cref{thm:p15-w02b-p15-073-invalid-candidate}(2) before any number is used.
The numbers exhibit the defect. In metres the candidate gives
\(1/(2+5)\) per second, about \(514\) per hour. In kilometres,
\(L=0.005\,\mathrm{km}\), it gives \(1/2.005\) per second, about \(1796\)
per hour. In millimetres, \(L=5000\,\mathrm{mm}\), it gives about \(0.72\)
per hour. The physical state is the same in all three, as part (3) of the
lemma predicts.
\item \(q\overset{?}{=}nv/\tau\). Here the claimed dimensionless product is
\(q\tau/(nv)\), with exponent vector \(a=(1,1,0,-1)\) over the columns
\(q,\tau,L,v\) of the matrix \(D\) of
\cref{sec:p15-w02b-p15-073-dimensions}. Then
\[
 Da=(0,-1)+(0,1)-(1,-1)=(-1,1)\neq0,
\]
so the product is not dimensionless: it carries the dimension of an inverse
speed. In the form of part (1), with inputs \(n,\tau,v\) and \(C=1\),
\([nv/\tau]=(1,-1)-(0,1)=(1,-2)\neq(0,-1)=[q]\), so the candidate is
invalid. In metres the candidate
gives \(15\) per second, or \(54\,000\) per hour; in kilometres,
\(v=0.03\,\mathrm{km/s}\), it gives \(54\) per hour.
\item \(q\overset{?}{=}n/(\tau+L/v)\). Now
\([L/v]=(1,0)-(1,-1)=(0,1)=[\tau]\), so the sum is valid, and
\([n/(\tau+L/v)]=(0,-1)=[q]\). The candidate passes both tests. In metres it
gives \(1/(2+1/6)\) per second, about \(1662\) per hour, and the value does
not change with the unit of length.
\end{enumerate}
Passing the check does not establish the third formula. The candidate
\(n/(\tau+2L/v)\) passes it equally. The coefficient of \(L/v\) comes from the
time-gap premise of \cref{def:p15-w02b-p15-073-lane-model}, through
\cref{thm:p15-w02b-p15-073-lane-flow}(2), not from dimensions. The first
candidate has a mechanical repair, dividing \(L\) by some speed, but the check
does not say which speed; that choice is again a model premise.
\end{hbexample}

\section{Results}
\label{sec:p15-w02b-p15-073-results}

\begin{hblemma}{thm:p15-w02b-p15-073-monotone-box}{proved=annex:p15-w02b-p15-073-monotone-box}{Ranges of monotone models and outward rounding}
Let \(f\) be continuous on a box \(B=\prod_i[l_i,u_i]\) and, in each
coordinate separately, either nondecreasing or nonincreasing.
Let \(\underline x\) be the corner taking \(l_i\) in each nondecreasing
coordinate and \(u_i\) in each nonincreasing one, and let \(\overline x\) be
the opposite corner. Then
\[
 f(B)=[f(\underline x),f(\overline x)].
\]
If \(m'\leq f(\underline x)\) and \(M'\geq f(\overline x)\), then
\(f(B)\subseteq[m',M']\). A threshold statement \(y>c\) is supported exactly
when \(f(\underline x)>c\).
\end{hblemma}

\begin{hbtheorem}{thm:p15-w02b-p15-073-lane-flow}{proved=annex:p15-w02b-p15-073-lane-flow}{Lane flow with and without the length term}
In steady single-lane traffic:
\begin{enumerate}
\item The car flow is \(v/s\), and the lane flow is \(q=nv/s\).
Consequently \(q\leq nv/L\) whatever the spacing rule.
\item In the time-gap model,
\[
 q=\frac{n}{\tau+L/v}=\frac{nv}{L+v\tau},\qquad
 \frac{q_0}{q}=1+\frac{L}{v\tau}=1+z .
\]
Hence \(q<q_0\) for every admissible input, and \(q\tau/n=1/(1+z)\).
\item On the box \(n\in[n_-,n_+]\), \(\tau\in[\tau_-,\tau_+]\),
\(L\in[L_-,L_+]\), \(v\in[v_-,v_+]\), with all endpoints positive,
the sets of values taken by \(q_0\) and by \(q\) are
\[
 \left[\frac{n_-}{\tau_+},\frac{n_+}{\tau_-}\right]
 \quad\text{and}\quad
 \left[\frac{n_-}{\tau_++L_+/v_-},\frac{n_+}{\tau_-+L_-/v_+}\right],
\]
respectively. The ratio \(q_0/q\) takes the values
\[
 \left[1+\frac{L_-}{v_+\tau_+},\;1+\frac{L_+}{v_-\tau_-}\right].
\]
\item If \(n\leq n_+\), \(L\geq L_-\) and
\(v<L_-n_-/(n_+\tau_+)\), then \(q<n_-/\tau_+\) under every spacing rule
with \(s\geq L\). Every such lane flow lies below the whole range of the
following-rule estimate on the box of part (3).
\item As \(v\to\infty\) with \(n,\tau,L\) fixed, \(q\to q_0\) and
\(q\sim q_0\); for each fixed \(v\) the bound \(q<q_0\) is strict.
\end{enumerate}
\end{hbtheorem}

Parts (1) and (4) use only the non-overlap premise \(s\geq L\). They therefore
survive any change in how drivers choose their gaps. Part (2) uses the
time-gap premise. Part (5) separates the asymptotic statement
\(q\sim q_0\) from a bound or identity, in the sense of
\cref{def:p15-approximation-relations}: the following-rule estimate is the
large-speed limit of the time-gap model, not its value at any finite speed.
Mahajan's sign \(\sim\) means something different: equality except perhaps for
a purely numerical factor of unstated size, with \(\approx\) reserved for a
factor near \(1\) and \(\propto\) for a factor that may carry dimensions
\cite[Sec.~1.1, p.~5]{Mahajan2014Insight}. The \(\sim\) of part (5) instead
asserts that the ratio \(q/q_0\) tends to exactly \(1\).

\section{Worked estimate}
\label{sec:p15-w02b-p15-073-worked}

\begin{hbexample}{def:p15-w02b-p15-073-lane-estimate}{One lane, with input ranges}
Stipulate the input box
\[
 1\leq n\leq1.5\ \text{people per car},\qquad
 1\,\mathrm s\leq\tau\leq2\,\mathrm s .
\]
The gap range runs from the source's daytime value to its compromise value;
the occupancy range widens the source's ``roughly one person'' to allow some
shared cars. Both ranges are stipulations for the example, not measurements.
Here and below times are in seconds, lengths in metres and speeds in metres
per second; a rate per second times \(3600\) is a rate per hour.
Since \(q_0=n/\tau\) is a monomial,
\cref{thm:p15-dimensional-product-bounds} gives its exact range. Multiplying
by \(3600\,\mathrm s\) per hour,
\[
 1800\leq q_0\leq5400\ \text{people per hour}.
\]
By \cref{def:p15-w02b-p15-073-reporting} the report is: \emph{between
\(1\cdot10^3\) and \(6\cdot10^3\) people per hour, of order \(10^3\)}.
The statement ``the lane carries at least \(10^3\) people per hour'' is
supported, since the lower endpoint \(1800\) exceeds \(10^3\)
(\cref{thm:p15-w02b-p15-073-monotone-box}).

The geometric midpoint of the range is about \(3118\). Reporting
``\(3118\) people per hour'' would claim four digits from a range that spans a
factor of three, which does not determine even the leading digit. The source's
\(1800\) is one corner of the box, the one with \(n=1\) and
\(\tau=2\,\mathrm s\), not a central value.
\end{hbexample}

\begin{hbexample}{def:p15-w02b-p15-073-highway}{The omitted length at highway speed}
Add the omitted quantities with stipulated ranges
\(4\,\mathrm m\leq L\leq6\,\mathrm m\) and
\(25\,\mathrm{m/s}\leq v\leq35\,\mathrm{m/s}\).
By \cref{thm:p15-w02b-p15-073-lane-flow}(3), the time-gap model gives
\[
 \frac{3600}{2+6/25}\approx1607\leq q\leq
 \frac{3600\cdot1.5}{1+4/35}\approx4846
\]
people per hour, and \(1+4/70\leq q_0/q\leq1+6/25=1.24\).
The rounded report is \emph{between \(1\cdot10^3\) and \(5\cdot10^3\) people
per hour, of order \(10^3\)}. The threshold of \(10^3\) is still supported,
since \(1607>10^3\). At highway speed the omission lowers the rounded upper
endpoint from \(6\cdot10^3\) to \(5\cdot10^3\) and leaves the order of
magnitude and the threshold verdict unchanged. Strictly, the following-rule
range does not contain every value of the fuller model, since
\(1607<1800\); but at each input point the two flows differ by a factor of
at most \(1.24\).
\end{hbexample}

\section{Where the estimate fails}
\label{sec:p15-w02b-p15-073-failure}

\begin{hbexample}{def:p15-w02b-p15-073-congested}{The same estimate in slow traffic}
Keep the ranges for \(n\), \(\tau\) and \(L\), and take
\(0.5\,\mathrm{m/s}\leq v\leq2\,\mathrm{m/s}\). By
\cref{thm:p15-w02b-p15-073-lane-flow}(3),
\[
 \frac{3600}{2+6/0.5}\approx257\leq q\leq\frac{3600\cdot1.5}{1+4/2}=1800
\]
people per hour, and \(2\leq q_0/q\leq13\).
The rounded report is \emph{between \(2\cdot10^2\) and \(2\cdot10^3\) people
per hour}, which does not lie within a single decade. Three conclusions of
the following-rule estimate fail:
\begin{enumerate}
\item The range \([1800,5400]\) excludes every value of \(q\) in this box
except the single corner value \(1800\), attained at \(n=1.5\),
\(\tau=1\,\mathrm s\), \(L=4\,\mathrm m\), \(v=2\,\mathrm{m/s}\).
\item The ratio \(q_0/q\) reaches \(13\), so the estimate can be wrong by
more than one order of magnitude.
\item The threshold statement \(q\geq10^3\) is no longer supported. Its
lower endpoint \(257\) lies below \(10^3\) and its upper endpoint \(1800\)
lies above it, so by \cref{thm:p15-approximation-certificates} the box
certifies neither \(q>10^3\) nor \(q<10^3\). The correct verdict is to abstain.
\end{enumerate}
These failures do not depend on the time-gap premise below the speed in
\cref{thm:p15-w02b-p15-073-lane-flow}(4): with \(L_-=4\,\mathrm m\),
\(n_-=1\), \(n_+=1.5\) and \(\tau_+=2\,\mathrm s\), every speed below
\(4/3\,\mathrm{m/s}\) gives \(q<1800\) people per hour under any spacing
rule that keeps cars from overlapping.
\end{hbexample}

The failure is located by the dimensionless group of
\cref{sec:p15-w02b-p15-073-dimensions}. The following-rule estimate is
adequate to a factor \(1+z\), with \(z=L/(v\tau)\). On the highway box
\(z\leq0.24\); on the slow box \(z\) reaches \(12\). Widening the ranges of
\(n\) and \(\tau\) cannot repair this, because no point of the original box
has \(z\) as an input. The repair is to add the omitted quantity to the model
and then to bound it.

\begin{hbwarning}{def:p15-w02b-p15-073-not-interval}{A range is not a confidence interval}
The interval \([1800,5400]\) is the deterministic image of a stipulated box.
No probability distribution was placed on \(n\) or \(\tau\), and the source's
gaps of \(1\) and \(4\) seconds describe different times of day, not sample
quantiles. No coverage probability attaches to the range. If a distribution
supported in the box were added, it would give the range probability one
under the model, and the slow-traffic case shows that the target can still
lie outside it: the probability is conditional on the model, and the model
omitted \(z\). Statistical uncertainty in measured inputs, rounding error and
model error are separate contributions, as in
\cref{def:p15-approximation-relations}.
\end{hbwarning}

\begin{hbmethod}{def:p15-w02b-p15-073-procedure}{Checking a range estimate for an omitted variable}
\begin{enumerate}
\item Write the target, the model \(f\), its listed inputs and their units.
Check that \(f\) has the dimension vector of the target and that every sum
in it adds terms of equal dimension vectors
(\cref{thm:p15-w02b-p15-073-invalid-candidate}); convert all inputs to one
stated system of units before substituting numbers.
Give each input a stipulated range and say what the range means.
\item Compute the exact image \(f(B)\), using monotonicity or
\cref{thm:p15-dimensional-product-bounds} where they apply, and report it by
\cref{def:p15-w02b-p15-073-reporting}.
\item State the conclusion to be drawn and test it against the endpoints.
\item List quantities that a fuller model would contain. Form the dimension
matrix of the target, the listed inputs and each candidate, and find the
additional dimensionless products it creates.
\item Find the range of each new product over the intended regime. Where it is
small, bound its effect on the conclusion; where it is not, supply the fuller
model or abstain from the conclusion.
\end{enumerate}
Side conditions: the box endpoints must be jointly admissible for the range to
be attained; restricting the box keeps the bounds but may lose attainment.
A product that is small in one regime need not be small in another.
\end{hbmethod}

\section{Limitations}
\label{sec:p15-w02b-p15-073-limits}

The time-gap model is a behavioural premise. In slow traffic drivers may keep
a distance rather than a time gap, lanes merge and change, and flow at a
point can be limited by a bottleneck downstream. Part (4) of
\cref{thm:p15-w02b-p15-073-lane-flow} avoids these by using only the
non-overlap condition, and so its conclusion holds for any of them; the
numerical ranges in the examples are stipulations and are not measured
traffic values. The same analysis does not settle Mahajan's rail--road
comparison, which also depends on the train model and on matching rush-hour
conditions. It shows only that a capacity figure read from the following rule
is an upper limit of the time-gap model, approached at high speed.

\section{Relationships and sources}
\label{sec:p15-w02b-p15-073-relations}

The rank count used for the omitted group is
\cref{thm:p15-dimensional-rank-nullity}, with the elimination procedure of
\cref{def:p15-dimensional-elimination}. The separation of the exponent
relation \(q\propto n/\tau\) at fixed \(z\) from the undetermined function
\(\Phi\) follows \cref{sec:p15-dimensional-pi}; the scaling family at fixed
\(z\) is the kind treated in \cref{def:p15-dimensional-scaling-law}.
Product-range bounds for monomial estimate trees are
\cref{thm:p15-dimensional-product-bounds}, and the source's other range
calculations are in \cref{sec:p15-dimensional-estimates}.
Decision certificates, factor bounds and the distinction between identity,
bound, asymptotic equivalence and heuristic estimate are in
\cref{sec:p15-approximate-quantitative-inference}.
Propagation of a probability model through a function, which a range
estimate does not supply, is treated in
\cref{sec:p15-measurement-uncertainty}.

The lane and train estimates, the two-second rule, the occupancy assumption,
and the qualification about varying gaps are from Mahajan's Section~1.2
\cite[pp.~6--7]{Mahajan2014Insight}; his table of graded equality signs is
in Section~1.1 \cite[p.~5]{Mahajan2014Insight}. The steady-traffic model, the
dimensional form, \cref{thm:p15-w02b-p15-073-lane-flow},
\cref{thm:p15-w02b-p15-073-monotone-box},
\cref{thm:p15-w02b-p15-073-invalid-candidate}, the candidate formulas of
\cref{def:p15-w02b-p15-073-invalid-check}, the reporting convention and the
numerical ranges are this handbook's derivations and stipulations; they are
not attributed to the source.
